% ===== §15 The gift: a threshold phenomenon across the three registers =====
\section{The Gift: A Threshold Phenomenon across the Three Registers}
\label{p09:sec:gift}

\noindent
There is a conspicuous absence in this paper, and its conspicuousness is the first thing to be said about it. We have built a theory of value across three registers---the imaginary, the symbolic, the real (\xref{p09:sec:value})---and a semiotics of the signs by which a relation is sustained or alienated (\xref{p09:sec:semiotics}); and yet we have passed over the one phenomenon in which all of these meet at once. \emb{The gift has been missing.} Its absence is not a mere oversight, a forgotten example to be supplied for completeness; it is structural, for the gift is precisely the phenomenon at which the three registers of value and the whole drama of the sign converge upon a single, weighted act. A paper about how value circulates across the three registers that did not, in the end, treat the gift would be a paper that had declined to test itself on its own hardest case.

The seventh paper of this series already marked the gift's peculiar density, and warned against the temptation it poses. There it was observed that the gift seems to sit naturally at the centre of every dimension at once---it has value, it is a sign, it must be read, it manufactures indebtedness, it transfers wealth across kin networks, and it is always given and received under the gaze of others---and that this very richness makes it tempting, and dangerous, to enthrone as the hub through which all else is read. \pinkref{Paper~VII} refused that enthronement, on the principle that load-bearing is conferred by a node's position in the relational network and not carried as its intrinsic essence. The present section honours that refusal: it treats the gift not as the hub of the theory but as its \emph{threshold phenomenon}---the place where the distinctions the paper has drawn are put to their sharpest test, precisely because the gift is the one act in which they are hardest to keep apart. In treating it here, after the apparatus is built rather than as the apparatus's organizing centre, the ninth paper and the seventh re-anchor one another once more.

\subsection{Why the gift is the threshold: the one materially three-register sign}
\label{p09:sec:gift-threshold}

What distinguishes the gift from the other signs the paper has treated---the declaration, the poem, the vow---is that the gift is \emph{materially} three-register at once. A poem, a declaration, operates mainly along the axis of the symbolic and the real; it has, as it were, no weight. The gift has weight. It cost money, or labour, or time, or some irreplaceable resource; it occupies the material world. And just because it is material, it is the superposition of all three registers of value in a single object:

\begin{claim}[The gift as three-register superposition]
The gift is the one sign that materially superposes all three registers of value. It is \emph{imaginary value}---the idealized token of the love it is taken to represent, the figure of ``what we are'' made into a thing. It is \emph{symbolic value}---countable, comparable, displayable, convertible: it has a price, a rank, a place in the economy of status and the ledger of who has given what to whom. And it is \emph{real value}---the trace, in a thing, of an investment that no price records: the unaccountable presence, the irreplaceable time, the witness of the other's unsayable need that the right gift somehow carries. The gift is thus the most concrete proving-ground of the theory of value: the question of which register \emph{leads} in a given gift is the question of what that gift, in truth, is.
\end{claim}

This superposition is why the gift is a threshold rather than a settled case. In the poem, the registers are relatively easy to keep apart; in the gift they are fused in a single object, and the same wrapped thing may be, for the giver, a witness of the real, and, for an onlooker, a move in the economy of status, and, for the receiver, the idealized token of a love that the giver does not in fact feel. The gift is where the registers are hardest to tell apart, and therefore where telling them apart matters most.

\subsection{The good gift: the poetic sign made material}
\label{p09:sec:good-gift}

Place the gift, now, against the criterion of the good sign (\xref{p09:sec:poem}). The true gift, I want to claim, is the poetic sign incarnate in matter---and the whole of Mauss's anthropology of the gift \parencite{mauss1990}, together with Hyde's account of the gift that must stay in motion \parencite{hyde1983}, read in this light, says so.

The decisive mark of the true gift is that \emb{it does not settle.} A settled exchange is a trade: I give you this, you give me its equivalent, and we are quits---the relation, having been balanced, is closed. The gift is the refusal of this closure. To return a gift with an exact equivalent, promptly, is not to honour the gift but to reject it: it declares the relation balanced, the account closed, and so refuses the very continuation the gift was offered to open. The gift obliges a return---this is Mauss's central finding \parencite{mauss1990}, the \emph{hau} or spirit of the thing given that compels reciprocity---but it obliges a return that must \emph{not} be equivalent, must come later, must itself be a gift and not a payment; and so the cycle of gift and counter-gift never closes, never settles, but spirals on, each gift carrying a part of the giver and opening rather than discharging the bond. \emb{This is the structure of the poem (\xref{p09:sec:poem}): the gift does not transmit a settled value but invites the receiver to walk a path of unprescribed response, and the value---the phase---is generated along that path.} The gift's meaning, like the poem's, is the holonomy accrued in the receiver's response; and because the response is free and unprescribed, each turn of the gift-cycle is a spiral and not a circle.

Read so, three traditions the paper has leaned on are revealed to be saying one thing at the gift. Mauss's gift that obliges an unequal, deferred return; Eglash's generative justice, in which value circulates back to its creators rather than being extracted to a centre; and the paper's own poetic cycle, in which the phase refluxes to every participant and is settled by none---these coincide at the true gift. \emb{The true gift is the material incarnation of the poetic, non-possessive cycle: it circulates value without extracting it, opens the bond without closing it, and returns to the giver not an equivalent but a continuation.}

\subsection{The vicious gift: settlement, bondage, and the siphon}
\label{p09:sec:vicious-gift}

But the gift is a threshold precisely because it falls so easily to the other side, and its corruptions track the three failures of the sign (\xref{p09:sec:failures}) with an exactness worth setting out.

\emc{The gift as settlement} is the gift degraded into disguised equivalent exchange: the present given to discharge a debt, to ``make up for'' an absence, to balance an account---the gift that says, beneath its wrapping, \emph{now we are quits}. This is settling failure (\xref{p09:sec:failures}) in material form: the use of a thing to cash out a relation, to declare closed what the true gift would have left open. The gift given to settle is not a gift but a payment that has not had the honesty to call itself one.

\emc{The gift as bondage} is the gift wielded to manufacture a dominating indebtedness: given not to open a free return but to bind, to place the receiver under an obligation they cannot discharge and so to subordinate them. Here the gift's non-closure, which in the true gift is the generous refusal to balance the account, is perverted into a debt deliberately kept open as a leash. This is the gift's specific way of reducing the other to fuel (\xref{p09:sec:fuel}): the relation is kept unbalanced not so that it may spiral but so that one party may hold the other in a permanent deficit.

\emc{The gift as siphon} is the gift in which symbolic value leads, and the object becomes a counter in the economy of status and the ledger of family capital---the competitive gift, the gift of face, the bride-price or dowry in its alienated form, in which what circulates is not the unaccountable real but the countable symbol, and what is secured is not a bond but a position. This is the gift assimilated to the symbolic register's liability to be siphoned (\xref{p09:sec:value}); and the alienation of social reproduction (\xref{p09:sec:coupled}) often runs through exactly such gifts, in which one party's uncounted real contribution is converted, in the family's ledger, into another's symbolic capital.

And over all three corruptions, the age of artificial generation casts its specific shadow. The gift that can be algorithmically recommended, mass-customized, generated at near-zero cost---the frictionless, optimized, dehumanized gift---is the material form of inflationary failure (\xref{p09:sec:failures}): the proliferation of the gift-signifier while the real investment of the signified goes to zero. A gift that cost the giver nothing of the unoutsourceable---no time, no presence, no attention that could not be delegated---is a gift drained of exactly the real value that made the gift a gift.

\subsection{Which register leads: the gift as diagnostic}
\label{p09:sec:gift-diagnostic}

The threshold character of the gift can now be stated as a diagnostic, and it gathers the section. Because the gift superposes all three registers, its goodness is determined by \emph{which register leads}---and the three possibilities map exactly onto the distinctions the paper has drawn.

\begin{claim}[The good gift is the one in which real value leads]
The \emph{imaginary-led} gift is the idealized token, the gift that figures a perfect image of the beloved or the bond---intense, and exhausting, for it encircles an idol (\xref{p09:sec:non-possess}) and erases the real other in favour of the image; it is the resplendent, brief gift, the rose made into an object. The \emph{symbolic-led} gift is the countable, displayable, comparable gift---the most readily settled, siphoned, and inflated; the gift of status and the ledger. The \emph{real-led} gift is the uncountable, void-encircling, witnessing gift---the one whose value lies not in the object but in the path of giving that was actually walked: the thing made by hand, the gift that cost the unoutsourceable, the present that witnesses the other's unsayable need without pretending to resolve it. Only the real-led gift bears the good's phase from within, and only it returns value to the giver-and-receiver without siphoning it elsewhere. The good gift is the one in which real value leads, and the object is merely the occasion of an encircling witness.
\end{claim}

So the gift, the threshold phenomenon, comes out exactly where the paper's whole argument would predict, and in coming out there confirms the argument from its hardest case. The same gift may be poem or contract, spiral or circle, the opening of a bond or the closing of an account, the return of value to its creators or the siphon that carries it off---and which it is depends not on the object, never on the object, but on whether it settles or invites, possesses or releases, encircles the idol or the void. The gift that leads with the real is the poetic cycle made into a thing one can hold; and that such a thing is possible---that value uncountable and unpossessable can be carried in an object that has a price and a weight---is perhaps the small everyday miracle on which the sustainability of intimacy, in its most concrete form, rests.
